%% ACM sigconf template — KDD conference format (double-column)
\documentclass[sigconf]{acmart}

\usepackage{booktabs}
\usepackage{multirow}
\usepackage{graphicx}
\usepackage{amsmath}
\usepackage{microtype}
\usepackage{xeCJK}
\setmainfont{Noto Serif CJK SC}
\setsansfont{Noto Sans CJK SC}
\setmonofont{Noto Sans Mono CJK SC}
\setCJKmainfont{Noto Serif CJK SC}
\setCJKsansfont{Noto Sans CJK SC}
\setCJKmonofont{Noto Sans Mono CJK SC}
\newcommand{\cn}[1]{#1}
\newcommand{\zhtext}[2]{\texorpdfstring{\cn{#1}}{#2}}

%% Remove ACM reference format / copyright box for course submission
\setcopyright{none}
\acmConference[CS306 DM Project]{Data Mining Course Final Project}{Spring 2026}{}
\acmDOI{}
\acmISBN{}

\begin{document}

\title{基于多维语言学特征工程的可解释 AIGC 检测：HC3 语料深度分析}

\author{\zhtext{杨意翀、韦嘉豪、熊佳漩、曾子致、江易明}{Yichong Yang, Jiahao Wei, Jiaxuan Xiong, Zizhi Zeng, Yiming Jiang}}
\affiliation{%
  \institution{\zhtext{南方科技大学}{Southern University of Science and Technology}}
}

%%----------------------------------------------------------
\begin{abstract}
随着 ChatGPT 等大语言模型（LLM）的快速普及，学术抄袭与虚假信息传播风险显著上升。
本文围绕“人类文本与 AI 生成文本（AIGC）如何区分”这一问题，基于 Human ChatGPT Comparison Corpus（HC3）
构建中英文双语可解释检测方案。我们不依赖黑盒深度模型，而是采用多维语言学特征工程：
从词汇丰富度、可读性、词性与句法、依存深度、话语连接与情感情绪六大维度提取每种语言 46--51 个手工特征。
通过皮尔逊相关性过滤（$|r|{>}0.85$）去除冗余后，训练并比较 Random Forest、XGBoost、LightGBM 三类集成树模型。
同时结合 SHAP 与多模型特征重要性共识解释模型决策。
实验结果表明，仅用手工特征即可在英文上达到 \textbf{95.1\%} Accuracy / F1\,=\,0.9506（AUC\,=\,0.9898），
在中文上达到 \textbf{92.1\%} Accuracy / F1\,=\,0.9199（AUC\,=\,0.9779）。
归因分析显示，二元词组重复、句法复杂度与词汇密度是最稳定、最具判别力的 AIGC 语言指纹。
\end{abstract}

\keywords{AIGC 检测，特征工程，可解释机器学习，SHAP，HC3，ChatGPT，文本挖掘}

\maketitle

%%----------------------------------------------------------
\section{引言}

Since its public release in November 2022, ChatGPT has demonstrated the ability
to produce fluent, coherent text across virtually any domain, challenging the
long-standing assumption that polished writing is exclusively a human skill.
While this capability unlocks tremendous productivity gains, it simultaneously
introduces pressing social risks: academic fraud, large-scale generation of
misinformation, and the erosion of trust in digital content~\cite{guo2023hc3}.

Existing AIGC detectors largely fall into two families.
\textit{Neural detectors} fine-tune pre-trained language models
(e.g., RoBERTa, DeBERTa) on labeled corpora and achieve near-ceiling accuracy
but are opaque, computationally expensive, and brittle to domain shift.
\textit{Zero-shot statistical methods} (e.g., DetectGPT~\cite{mitchell2023detectgpt},
GLTR~\cite{gehrmann2019gltr}) exploit log-probability perturbations from the generating
model itself, but require white-box access to the LLM.

We argue that neither family adequately serves the practitioner who needs a
detector that is (a)~\emph{accurate}, (b)~\emph{interpretable}, and
(c)~\emph{model-agnostic} (applicable without access to the suspected generator).
Our response is a feature-engineering-first approach that systematically mines the
statistical and linguistic regularities that separate human from AI writing.

\paragraph{主要贡献。}
\begin{enumerate}
\item We construct a rich, \emph{bilingual} (English + Chinese) feature space
      of 46--51 linguistic features organized into six interpretable categories.
\item We demonstrate that gradient-boosted tree classifiers trained on this feature
      space achieve 92--95\% accuracy on the HC3 benchmark, rivaling fine-tuned
      neural methods while remaining fully transparent.
\item We provide the first SHAP-based attribution study on the HC3 corpus that
      quantitatively ranks which specific feature combinations expose AI text,
      and we validate the finding with multi-model consensus across three classifiers.
\end{enumerate}

%%----------------------------------------------------------
\section{相关工作}

\paragraph{HC3 与语言学分析。}
Guo et al.~\cite{guo2023hc3} introduced the HC3 dataset and conducted linguistic
comparisons showing that ChatGPT answers tend to be longer, more neutral in sentiment,
and richer in discourse connectives.
Their fine-tuned RoBERTa detector reached $\sim$98.8\% F1 on English
but offers no interpretability.

\paragraph{统计/零样本检测方法。}
GLTR~\cite{gehrmann2019gltr} visualises per-token rank distributions from a
reference LM to expose machine text.
DetectGPT~\cite{mitchell2023detectgpt} perturbs candidate texts and measures
log-probability curvature; it needs no training data but requires LLM access.

\paragraph{文体计量与特征工程检测。}
Stylometric approaches (Type-Token Ratio, Honore's R, readability scores)
have been applied to authorship attribution for decades~\cite{juola2006}.
Recent work has adapted these to AIGC detection~\cite{bhatt2023stylometric},
but most studies focus on English only and do not provide SHAP attribution.
Our work extends these directions to a bilingual setting with a comprehensive
SHAP analysis.

%%----------------------------------------------------------
\section{数据集}

\subsection{Human ChatGPT Comparison Corpus（HC3）}

We use the HC3 dataset~\cite{guo2023hc3} as our \emph{single, deep-analysed} corpus,
covering both English (\textbf{HC3-EN}) and Chinese (\textbf{HC3-ZH}) subsets.

\paragraph{HC3-EN} collects 24\,322 questions with a total of 58\,546 human answers
(sourced from Reddit ELI5, WikiQA, Wikipedia, medical dialogues, and financial QA)
and 26\,903 ChatGPT answers across five domain splits:
\textit{reddit\_eli5}, \textit{open\_qa}, \textit{wiki\_csai}, \textit{medicine}, and
\textit{finance}.

\paragraph{HC3-ZH} collects 12\,853 questions with 22\,259 human answers and
17\,522 ChatGPT answers across seven domain splits:
\textit{open\_qa}, \textit{baike}, \textit{nlpcc\_dbqa}, \textit{medicine},
\textit{finance}, \textit{psychology}, and \textit{law}.

\subsection{预处理与采样}

All answers shorter than 30 characters (English) or 20 characters (Chinese) are
discarded to remove near-empty responses.
To keep runtime tractable and class balance exact,
we draw a stratified random sample of 10\,000 texts (5\,000 human, 5\,000 AI)
from each language using a fixed seed (42).
The final splits are summarized in Table~\ref{tab:dataset}.

\begin{table}[h]
\centering
\caption{预处理与分层采样后的数据集统计。}
\label{tab:dataset}
\small
\begin{tabular}{lcccc}
\toprule
\textbf{Split} & \textbf{Lang.} & \textbf{Human} & \textbf{AI} & \textbf{Total} \\
\midrule
Train (80\%) & EN & 4000 & 4000 & 8000 \\
Test  (20\%) & EN & 1000 & 1000 & 2000 \\
Train (80\%) & ZH & 4000 & 4000 & 8000 \\
Test  (20\%) & ZH & 1000 & 1000 & 2000 \\
\bottomrule
\end{tabular}
\end{table}

%%----------------------------------------------------------
\section{特征工程}

Our feature space is organized into six interpretable categories.
English texts yield 51 raw features; Chinese texts yield 46, reflecting the
language-specific adaptations described below.

\subsection{类别一：词汇丰富度}

Lexical diversity measures the breadth of vocabulary employed by a writer.

\begin{itemize}
\item \textbf{Type-Token Ratio (TTR)}: $\text{TTR} = V / N$ where $V$ is the number
  of unique tokens and $N$ the total token count.
  AI text tends to have lower TTR due to formulaic phrasing.
\item \textbf{Hapax Legomena Ratio}: the fraction of words appearing exactly once,
  reflecting authentic creative choices unique to human writers.
\item \textbf{Honore's R}: $R = 100 \cdot \log N / (1 - V_1/V)$ where $V_1$ is the
  hapax count. This statistic penalises vocabulary reuse more heavily than TTR.
\item \textbf{Vocabulary Density}: $V / \sqrt{N}$, a length-normalised diversity score.
\end{itemize}

For Chinese, word-level (\texttt{ttr\_word}) and character-level (\texttt{ttr\_char})
TTR are computed separately because Chinese word segmentation is ambiguous.

\subsection{类别二：可读性}

For English we use four scores from the \texttt{textstat} library:
Flesch Reading Ease, Flesch-Kincaid Grade Level, average word length (characters),
and average syllables per word.
For Chinese, which lacks syllabic readability formulae, we substitute proxies:
average word length in characters, average sentence length in characters, and
average sentence length in words.

\subsection{类别三：词性与句法特征}

Part-of-speech tag ratios capture grammatical style.
Eight ratios are computed for English (noun, verb, adjective, adverb, conjunction,
determiner, punctuation, pronoun) via NLTK's averaged perceptron tagger.
Chinese POS tagging uses jieba's hidden Markov model tagger with equivalent
categories (noun \textit{n}/\textit{nr}/\textit{ns}/\textit{nt}/\textit{nz},
verb \textit{v}/\textit{vd}/\textit{vn}, etc.).

Additionally, for English only, we parse each text with spaCy's dependency parser
and compute five \textbf{dependency tree depth} features:
maximum depth, average depth, standard deviation of depth,
ratio of deep clauses (depth~$>$~5), and average branching factor.
These capture the hierarchical complexity of clause nesting, which is notably
higher in AI-generated English.

\subsection{类别四：话语结构（连接词）}

AI text is hypothesised to overuse structured discourse markers.
We define 7 sub-categories of transition words (Table~\ref{tab:transitions})
and compute per-category frequency normalized per 100 words.
An aggregate density and a \textit{structured marker ratio}
(sequential + conclusive markers) are also computed.

\begin{table}[h]
\centering
\caption{连接词子类别与示例。}
\label{tab:transitions}
\small
\begin{tabular}{ll}
\toprule
\textbf{Category} & \textbf{Examples (EN)} \\
\midrule
Sequential   & firstly, secondly, finally \\
Conclusive   & in conclusion, to summarize, overall \\
Contrastive  & however, nevertheless, on the other hand \\
Additive     & moreover, furthermore, in addition \\
Causal       & therefore, consequently, thus \\
Exemplifying & for example, for instance, such as \\
Hedging      & it is important to note, generally speaking \\
\bottomrule
\end{tabular}
\end{table}

\subsection{类别五：情感与情绪}

For English, VADER sentiment analysis yields compound, neutral, negative,
and positive scores.
An \textit{emotion word ratio} is computed using a fixed lexicon of 18 emotionally
charged words.
Three \textit{per-sentence} statistics quantify emotional variability across the text:
fluctuation (standard deviation of per-sentence compound scores),
range (max~$-$~min), and \textit{flip ratio} (fraction of consecutive sentences where
sentiment polarity changes).
The intuition is that human writing exhibits more affective variation—the
``machine flavor'' of AI text is its flat, consistently neutral register.

For Chinese, VADER is unavailable; we instead use a curated positive/negative
lexicon (20+18 words) and compute analogous ratio and fluctuation features,
supplemented by exclamation and question mark ratios.

\subsection{类别六：文本统计特征}

Length (characters, words), sentence count, average sentence length,
stopword ratio, capitalization ratio (English), unique bigram ratio, digit ratio.
For Chinese, an additional \textit{four-character phrase ratio} captures usage of
chengyu (classical idiom phrases), which appear more often in formal human writing.

\subsection{相关性过滤}

To mitigate multicollinearity and reduce feature dimensionality,
we compute the pairwise Pearson correlation matrix on the training set and
remove one feature from each pair with $|r|{>}0.85$,
retaining the member with higher variance.
This step reduces the English feature space from 51 to \textbf{42} features
and the Chinese space from 46 to \textbf{40} features (Table~\ref{tab:removed}).

\begin{table}[h]
\centering
\caption{相关性过滤移除特征（$|r|{>}0.85$）。}
\label{tab:removed}
\small
\begin{tabular}{ll}
\toprule
\textbf{Language} & \textbf{Removed features} \\
\midrule
English (9) & \texttt{sentence\_count}, \texttt{ttr}, \texttt{avg\_syllables\_per\_word}, \\
            & \texttt{flesch\_kincaid\_grade}, \texttt{avg\_sent\_len}, \\
            & \texttt{sentiment\_fluctuation}, \texttt{trans\_sequential}, \\
            & \texttt{text\_len\_words}, \texttt{avg\_para\_len} \\
\midrule
Chinese (6) & \texttt{avg\_sent\_len}, \texttt{sentiment\_fluctuation}, \\
            & \texttt{trans\_sequential}, \texttt{text\_len\_words}, \\
            & \texttt{avg\_sent\_len\_words}, \texttt{ttr\_word} \\
\bottomrule
\end{tabular}
\end{table}

%%----------------------------------------------------------
\section{实验设置}

\subsection{模型设置}

Three ensemble tree classifiers are trained:

\textbf{Random Forest (RF):}
200 trees, max depth 15, minimum samples per leaf 10, trained in parallel.

\textbf{XGBoost:}
200 rounds, max depth 6, learning rate 0.1, subsample and column-subsample 0.8,
log-loss evaluation metric.

\textbf{LightGBM:}
200 rounds, max depth 8, learning rate 0.1, subsample and column-subsample 0.8.

All features are standardized with zero mean and unit variance (\texttt{StandardScaler})
before training.
An 80/20 stratified train-test split (random state 42) is used for all experiments.
Additionally, we report 5-fold stratified cross-validation F1 on the training set
using LightGBM (100 rounds) to assess stability.

\subsection{评价指标}

We report Accuracy, Precision, Recall, F1-Score (macro over both classes), and
ROC-AUC.

%%----------------------------------------------------------
\section{实验结果}

\subsection{HC3 英文子集}

Table~\ref{tab:en_results} reports test-set performance for all three models.
XGBoost achieves the best overall result: \textbf{95.1\%} accuracy,
F1~=~0.9506, AUC~=~0.9898.
LightGBM is a close second with a 5-fold CV mean F1 of 0.9394 (±0.0040),
confirming the result is stable and not a lucky train-test split.

\begin{table}[h]
\centering
\caption{HC3-EN 测试集结果（2000 样本）。}
\label{tab:en_results}
\small
\begin{tabular}{lccccc}
\toprule
\textbf{Model} & \textbf{Acc.} & \textbf{Prec.} & \textbf{Rec.} & \textbf{F1} & \textbf{AUC} \\
\midrule
Random Forest & 0.9300 & 0.9479 & 0.9100 & 0.9286 & 0.9798 \\
XGBoost       & \textbf{0.9510} & \textbf{0.9583} & 0.9430 & \textbf{0.9506} & \textbf{0.9898} \\
LightGBM      & 0.9490 & 0.9554 & \textbf{0.9420} & 0.9486 & 0.9887 \\
\midrule
\multicolumn{6}{l}{\textit{5-Fold CV (LightGBM): mean F1 = 0.9394 ± 0.0040}} \\
\bottomrule
\end{tabular}
\end{table}

\subsection{HC3 中文子集}

Results on Chinese (Table~\ref{tab:zh_results}) are slightly lower, reflecting the
harder linguistic disambiguation problem. LightGBM edges XGBoost:
\textbf{92.1\%} accuracy, F1~=~0.9199, AUC~=~0.9779.
The 5-fold CV mean F1 of 0.9165 (±0.0056) is equally stable.

\begin{table}[h]
\centering
\caption{HC3-ZH 测试集结果（2000 样本）。}
\label{tab:zh_results}
\small
\begin{tabular}{lccccc}
\toprule
\textbf{Model} & \textbf{Acc.} & \textbf{Prec.} & \textbf{Rec.} & \textbf{F1} & \textbf{AUC} \\
\midrule
Random Forest & 0.9000 & 0.9049 & 0.8940 & 0.8994 & 0.9665 \\
XGBoost       & 0.9200 & 0.9242 & 0.9150 & 0.9196 & 0.9782 \\
LightGBM      & \textbf{0.9205} & \textbf{0.9269} & 0.9130 & \textbf{0.9199} & 0.9779 \\
\midrule
\multicolumn{6}{l}{\textit{5-Fold CV (LightGBM): mean F1 = 0.9165 ± 0.0056}} \\
\bottomrule
\end{tabular}
\end{table}

\subsection{跨语言对比}

English is harder to fool the detector in: the XGBoost model reaches 95.1\% vs.\
92.1\% for Chinese. One likely factor is the availability of dependency parse features
(spaCy, EN only) which contribute 13.3\% of total SHAP importance.
Chinese compensates with richer POS taxonomy (jieba) and character-level features
unavailable in English.

%%----------------------------------------------------------
\section{SHAP 归因分析}

\subsection{SHAP 重要特征排序}

We apply TreeExplainer~\cite{lundberg2017shap} to the XGBoost model on a 500-sample
subset of the test set. Table~\ref{tab:shap_en} and Table~\ref{tab:shap_zh} list the
top-10 features ranked by mean $|\text{SHAP}|$ value, together with the direction of
their effect on the ``AI'' class probability.

\begin{table}[h]
\centering
\caption{HC3-EN（XGBoost）SHAP Top-10 特征。}
\label{tab:shap_en}
\small
\begin{tabular}{clcc}
\toprule
\textbf{Rank} & \textbf{Feature} & \textbf{SHAP} & \textbf{Direction} \\
\midrule
1  & \texttt{unique\_bigram\_ratio}  & 1.4193 & $\uparrow$ in AI \\
2  & \texttt{dep\_depth\_avg}        & 0.7910 & $\uparrow$ in AI \\
3  & \texttt{hapax\_ratio}           & 0.6410 & $\downarrow$ in AI \\
4  & \texttt{adv\_ratio}             & 0.6175 & $\downarrow$ in AI \\
5  & \texttt{vocab\_density}         & 0.5275 & $\downarrow$ in AI \\
6  & \texttt{text\_len\_chars}       & 0.5098 & $\downarrow$ in AI \\
7  & \texttt{stopword\_ratio}        & 0.4575 & $\downarrow$ in AI \\
8  & \texttt{sent\_len\_cv}          & 0.4020 & $\downarrow$ in AI \\
9  & \texttt{punct\_ratio}           & 0.3110 & $\downarrow$ in AI \\
10 & \texttt{paragraph\_count}       & 0.2871 & $\uparrow$ in AI \\
\bottomrule
\end{tabular}
\end{table}

\begin{table}[h]
\centering
\caption{HC3-ZH（XGBoost）SHAP Top-10 特征。}
\label{tab:shap_zh}
\small
\begin{tabular}{clcc}
\toprule
\textbf{Rank} & \textbf{Feature} & \textbf{SHAP} & \textbf{Direction} \\
\midrule
1  & \texttt{unique\_bigram\_ratio}  & 0.7501 & $\uparrow$ in AI \\
2  & \texttt{vocab\_density}         & 0.6175 & $\downarrow$ in AI \\
3  & \texttt{adv\_ratio\_zh}         & 0.6094 & $\downarrow$ in AI \\
4  & \texttt{pronoun\_ratio\_zh}     & 0.5669 & $\uparrow$ in AI \\
5  & \texttt{sent\_len\_cv}          & 0.5117 & $\downarrow$ in AI \\
6  & \texttt{conj\_ratio\_zh}        & 0.4234 & $\uparrow$ in AI \\
7  & \texttt{honore\_r}              & 0.4022 & $\downarrow$ in AI \\
8  & \texttt{trans\_hedging}         & 0.3764 & $\uparrow$ in AI \\
9  & \texttt{avg\_para\_len}         & 0.3322 & $\uparrow$ in AI \\
10 & \texttt{avg\_word\_len\_chars}  & 0.3258 & $\downarrow$ in AI \\
\bottomrule
\end{tabular}
\end{table}

\subsection{特征组重要性}

Table~\ref{tab:group} aggregates SHAP values by the six feature categories.
For English, \textit{Text Statistics} dominates (26.6\%), followed by
\textit{POS \& Syntax} (18.1\%) and \textit{Dependency Depth} (13.3\%).
For Chinese, \textit{POS \& Syntax} leads (26.7\%), reflecting the greater
diagnostic power of the jieba POS taxonomy.

\begin{table}[h]
\centering
\caption{特征组 SHAP 贡献占比。}
\label{tab:group}
\small
\begin{tabular}{lcc}
\toprule
\textbf{Feature Group} & \textbf{EN (\%)} & \textbf{ZH (\%)} \\
\midrule
Text Statistics        & 26.6 & 19.1 \\
POS \& Syntax          & 18.1 & 26.7 \\
Dependency Depth       & 13.3 & ---  \\
Vocabulary Richness    & 11.3 & 15.9 \\
Sentence Structure     & 9.8  & 18.0 \\
Transition Words       & 8.8  & 8.5  \\
Sentiment \& Emotion   & 6.9  & 5.0  \\
Readability            & 5.2  & 6.8  \\
\bottomrule
\end{tabular}
\end{table}

\subsection{多模型特征重要性共识}

To verify that the SHAP ranking is not model-specific, we rank features by impurity-based
importance in all three models and compute a consensus rank (sum of ranks; lower is better).
Table~\ref{tab:consensus_en} shows the top-7 consensus features for English.
All three models unanimously rank \texttt{unique\_bigram\_ratio} first (★★★),
strongly validating its diagnostic power.

\begin{table}[h]
\centering
\caption{多模型共识（HC3-EN）。★★★ 表示最高共识等级。}
\label{tab:consensus_en}
\small
\begin{tabular}{lcccc}
\toprule
\textbf{Feature} & \textbf{RF} & \textbf{XGB} & \textbf{LGB} & \textbf{Tier} \\
\midrule
\texttt{unique\_bigram\_ratio} & 1 & 1 & 1 & ★★★ \\
\texttt{dep\_depth\_avg}       & 3 & 2 & 4 & ★★★ \\
\texttt{text\_len\_chars}      & 5 & 6 & 3 & ★★★ \\
\texttt{hapax\_ratio}          & 2 & 3 & 9 & ★★★ \\
\texttt{vocab\_density}        & 7 & 8 & 6 & ★★ \\
\texttt{sent\_len\_cv}         & 10& 9 & 2 & ★★ \\
\texttt{adv\_ratio}            & 8 & 12& 11& ★★ \\
\bottomrule
\end{tabular}
\end{table}

%%----------------------------------------------------------
\section{关键发现与讨论}

\paragraph{发现 1：二元词组重复是通用 AI 指纹。}
\texttt{unique\_bigram\_ratio} (higher in AI) ranks first by SHAP in both languages
and is the only ★★★ consensus feature in both English and Chinese models.
ChatGPT generates text with more repetitive, formulaic word-pair combinations
than human writers across all domains.

\paragraph{发现 2：英文中句法复杂度具有显著区分力。}
Average dependency tree depth (\texttt{dep\_depth\_avg}) is the \#2 feature in English,
contributing 13.3\% of group SHAP. AI-generated English text features significantly
deeper clause hierarchies—subordinate clauses nested within subordinate clauses—a
pattern consistent with ChatGPT's instruction-following optimization for completeness.

\paragraph{发现 3：中文检测主要由词性分布驱动。}
In Chinese, POS \& Syntax features collectively account for 26.7\% of total SHAP.
Humans use more adverbs (\texttt{adv\_ratio\_zh} $\downarrow$ in AI), reflecting
colloquial hedging and emotional intensification, while AI overuses conjunctions
(\texttt{conj\_ratio\_zh} $\uparrow$ in AI) to enforce explicit logical structure.

\paragraph{发现 4：AI 文本词汇多样性持续偏低。}
Both \texttt{hapax\_ratio} and \texttt{vocab\_density} are negatively associated with
AI text across both languages. Despite ChatGPT's vast training vocabulary, its
output reuses a more constrained lexical set within any given response, possibly due
to next-token prediction converging to high-probability, domain-common words.

\paragraph{发现 5：人类文本句长波动更大。}
\texttt{sent\_len\_cv} (coefficient of variation of sentence lengths) is lower in AI text
in both languages, indicating that ChatGPT tends to produce sentences of more uniform
length—another manifestation of its ``machine-flavored'' regularity.

\paragraph{发现 6：连接与话语特征验证 AI 结构化表达倾向。}
Hedging expressions (\texttt{trans\_hedging}) rank \#8 in Chinese SHAP, confirming
that phrases like ``\zhtext{值得注意的是}{值得注意的是}'' (``it is worth noting'') and ``\zhtext{一般来说}{一般来说}''
(``generally speaking'') are disproportionately frequent in AI text.
Collectively, transition word features contribute 8.5--8.8\% of total SHAP,
a modest but consistent contribution.

\paragraph{情绪平坦性：虽被过滤但有效的信号。}
\texttt{sentiment\_fluctuation} was removed by correlation filtering in both languages
because it correlates ($|r|{>}0.85$) with other sentiment statistics.
However, \texttt{sentiment\_neutrality} (higher in AI) survives and contributes to the
Sentiment \& Emotion group. This supports the hypothesis that AI text lacks genuine
emotional variation—the well-documented ``machine flavor'' of tonally flat prose.

\paragraph{与神经网络基线对比。}
Our 92--95\% accuracy using only hand-crafted features compares favorably to
the HC3 baseline detectors: Guo et al.\ report $\sim$98.8\% F1 with fine-tuned
RoBERTa~\cite{guo2023hc3}, a 3--7\% gap that is the price of full interpretability.
Unlike RoBERTa, our model explains every prediction through SHAP values,
identifies the specific linguistic features driving each decision, and requires
no GPU infrastructure at inference time.

%%----------------------------------------------------------
\section{结论}

We presented an interpretable, feature-engineering-first pipeline for bilingual
AIGC detection using the HC3 corpus.
Six categories of 40--42 hand-crafted linguistic features—spanning vocabulary
richness, readability, syntactic structure, discourse patterns, and sentiment—feed
three gradient-boosted tree classifiers (RF, XGBoost, LightGBM).
SHAP attribution quantitatively reveals that \emph{bigram repetition, syntactic
clause depth, and vocabulary diversity} are the dominant signals exposing AI text,
with consistent cross-language and cross-model agreement.

Our work demonstrates that accurate, interpretable AIGC detection is achievable
without black-box deep learning. The feature taxonomy and SHAP methodology developed
here can readily be extended to new languages, new LLMs, and domain-specific
detection scenarios.

\paragraph{局限性与未来工作。}
The current feature set does not model \emph{topic or domain distribution} shifts;
future work should evaluate cross-domain generalization.
Dependency parsing is available for English only; adding a Chinese dependency parser
(e.g., LTP) could narrow the 3\% accuracy gap.
As LLMs evolve, the feature thresholds may drift; periodic recalibration will be needed.

%%----------------------------------------------------------
\begin{acks}
感谢 HC3 数据集贡献者（Guo 等）公开语料，为 AIGC 检测提供了可复现研究基础。
\end{acks}

%%----------------------------------------------------------
\renewcommand{\refname}{参考文献}
\renewcommand{\bibname}{参考文献}
\bibliographystyle{ACM-Reference-Format}
\begin{thebibliography}{9}

\bibitem{guo2023hc3}
Biyang Guo, Xin Zhang, Ziyuan Wang, Minqi Jiang, Jinran Nie, Yuxuan Ding, Jianwei Yue, and Yupeng Wu.
\newblock How Close is ChatGPT to Human Experts? Comparison Corpus, Evaluation, and Detection.
\newblock \textit{arXiv 预印本 arXiv:2301.07597}, 2023。

\bibitem{mitchell2023detectgpt}
Eric Mitchell, Yoonho Lee, Alexander Khazatsky, Christopher D.\ Manning, and Chelsea Finn.
\newblock DetectGPT: Zero-Shot Machine-Generated Text Detection using Probability Curvature.
\newblock 载于 \textit{ICML 会议论文集}, 2023。

\bibitem{gehrmann2019gltr}
Sebastian Gehrmann, Hendrik Strobelt, and Alexander M.\ Rush.
\newblock GLTR: Statistical Detection and Visualization of Generated Text.
\newblock 载于 \textit{ACL Demo Track 论文集}, 2019。

\bibitem{lundberg2017shap}
Scott M.\ Lundberg and Su-In Lee.
\newblock A Unified Approach to Interpreting Model Predictions.
\newblock 载于 \textit{NeurIPS 会议论文集}, 2017。

\bibitem{juola2006}
Patrick Juola.
\newblock Authorship Attribution.
\newblock \textit{Foundations and Trends in Information Retrieval}, 1(3):233--334, 2006。

\bibitem{bhatt2023stylometric}
Gaurav Bhatt, Bhavna Jain, and Sanjay Chaudhary.
\newblock Stylometric Detection of AI-Generated Text in Twitter Timelines.
\newblock \textit{arXiv 预印本 arXiv:2303.03697}, 2023。

\end{thebibliography}

\end{document}
